\documentclass{article}

% Pass the round citation option to natbib BEFORE the NeurIPS style auto-loads it.
% The style loads natbib itself, so a second \usepackage[round]{natbib} would cause
% an "Option clash for package natbib" error; passing the option here avoids it.
\PassOptionsToPackage{round}{natbib}

% Keep the same paper style when the NeurIPS style file is available.
\IfFileExists{neurips_2026.sty}{\usepackage[preprint]{neurips_2026}}{}

\usepackage[utf8]{inputenc}
\usepackage[T1]{fontenc}
\usepackage{url}
\usepackage{booktabs}
\usepackage{amsfonts}
\usepackage{microtype}
\usepackage{amsmath}
\usepackage{amssymb}
\usepackage{graphicx}
\usepackage{float}
\usepackage{enumitem}
% natbib is loaded by the NeurIPS style above (with the round option passed via
% \PassOptionsToPackage near \documentclass); do not load it again here.
\usepackage[hidelinks,colorlinks=true,citecolor=blue,linkcolor=blue,urlcolor=blue]{hyperref}
\usepackage[capitalize,noabbrev]{cleveref}

% >>> BEGIN COPY INTO A DOCUMENT PREAMBLE >>>
% Copied verbatim from the shared skill's style reference
% /p/rlprojects/.claude/skills/rnd-experiment-tex-track/example.tex (its marked block).
% WARNING: \arrayrulecolor is a GLOBAL switch — only use it through the two rule macros below,
% which both end with the color on black. A \multicolumn that is NOT the last cell of its row
% must end its spec with ? (e.g. {c?}), or the grey vertical on its right edge disappears.
\usepackage[table]{xcolor}  % [table] loads colortbl -> \arrayrulecolor
\usepackage{array}
\setlength{\arrayrulewidth}{0.5pt}                 % \hline/\vline only; booktabs rules unaffected
\definecolor{tabgrey}{gray}{0.65}
\newcommand{\greyvline}{{\color{tabgrey}\vline}}   % braces keep the color inside the rule
\newcolumntype{?}{!{\greyvline}}                   % column separator: one grey vertical rule
\newcommand{\greyrule}{\arrayrulecolor{tabgrey}\hline\arrayrulecolor{black}} % between rows of one category
\newcommand{\blackrule}{\arrayrulecolor{black}\hline}                        % category boundary / frame
% <<< END COPY <<<

% Read / unread / updated marking for the contents table (scripts/report_contents.py writes the
% \study...begin and \study...end lines, and a \studyfloatcolor line at the start of every table and
% figure of a marked subsection, because a float does not keep the text's color).
\definecolor{unreadblue}{RGB}{0,70,200}
\definecolor{editbrown}{RGB}{133,45,15}
\newif\ifstudyunread
\newif\ifstudychanged
\newcommand{\studyunreadbegin}{\global\studyunreadtrue\color{unreadblue}}
\newcommand{\studyunreadend}{\global\studyunreadfalse\color{black}}
\newcommand{\studychangedbegin}{\global\studychangedtrue\color{editbrown}}
\newcommand{\studychangedend}{\global\studychangedfalse\color{black}}
\newcommand{\studyfloatcolor}{\ifstudyunread\color{unreadblue}\fi\ifstudychanged\color{editbrown}\fi}

% The NeurIPS style sets \flushbottom, which stretches a page's vertical spacing until the text
% reaches the bottom margin. This document's pages are mostly one tall table each, so a page that
% a table did not fit on gets its few paragraphs spread across the whole height with large gaps
% between them. \raggedbottom leaves that space at the bottom instead.
\raggedbottom

\title{Exploration Bonuses on the JAX Training Platform:\\ Environment and Train-Run Development Document}

\author{%
  Anonymous Author(s)\\
  Affiliation\\
  \texttt{email}
}

% Tables and figures are numbered by section (Table 1.1, Figure 1.2, ...), so a reader who is
% told "the table in 1.1" can find it, and so a later section's artifacts never renumber this
% section's. amsmath, loaded above, provides \numberwithin.
\numberwithin{table}{section}
\numberwithin{figure}{section}

\begin{document}

\maketitle

\tableofcontents
% >>> REPORT CONTENTS START
\begin{center}
\small
\begin{tabular}{@{}>{\raggedright\arraybackslash}p{6.4cm} l l l@{}}
\toprule
\textbf{subsection} & \textbf{first written} & \textbf{last changed} & \textbf{status} \\
\midrule
\hyperref[sec:pm-jax-envspec]{Environment specification} & 2026-08-16 19:06 PT & 2026-10-13 20:24 PT & read \\
\hyperref[sec:pm-jax-run11]{Train run 1.1: the first algorithm batch} & 2026-08-16 19:06 PT & 2026-10-19 12:58 PT & read \\
\hyperref[sec:pm-jax-corrected]{Train run 1.2: corrected intrinsic returns} & 2026-10-13 22:23 PT & 2026-10-19 12:58 PT & read \\
\textcolor{unreadblue}{\hyperref[sec:pm-jax-speed-jaguar03]{Training speed on jaguar03}} & 2026-10-19 16:44 PT & 2026-10-19 16:44 PT & unread \\
\textcolor{unreadblue}{\hyperref[sec:pm-jax-screening]{Choosing a small environment for the reward-decline study}} & 2026-10-19 16:55 PT & 2026-10-19 16:55 PT & unread \\
\textcolor{unreadblue}{\hyperref[sec:pm-jax-decline-batch1]{Reward-decline batch 1: a control and ten ideas}} & 2026-10-19 16:55 PT & 2026-10-19 16:55 PT & unread \\
\bottomrule
\end{tabular}
\end{center}
{\small Times come from the document's version history and are Pacific, converted from machines running Eastern. Unread subsections are \textcolor{unreadblue}{blue}; in updated ones only the changed text is \textcolor{editbrown}{brown}.\par}
% <<< REPORT CONTENTS END
\clearpage

%==============================================================================
\section{PointMaze on the JAX platform}
\label{sec:pm-jax}
% The environment is subsection 1.0 and the first train run is 1.1, so a later run on the same
% environment becomes 1.2 and the shared environment specification keeps the .0 slot.
\setcounter{subsection}{-1}

\subsection{Environment specification}
\label{sec:pm-jax-envspec}

The platform uses the large PointMaze with fixed start and goal positions and no reset noise.
\Cref{tab:pm-jax-env} gives the environment specification;
\Cref{fig:pm-jax-environment} shows the maze and the grids used by the visit-count bonuses.

\begin{table}[H]
\centering
\footnotesize
\setlength{\tabcolsep}{3pt}
\renewcommand{\arraystretch}{1.2}
\begin{tabular}{@{}>{\raggedright\arraybackslash}p{3.4cm}
?>{\raggedright\arraybackslash}p{\dimexpr\textwidth-3.4cm-2\tabcolsep-\arrayrulewidth\relax}@{}}
\blackrule
\textbf{Parameter} & \textbf{PointMaze large, no reset noise} \\
 & \small(the platform's only environment so far) \\
\blackrule
\multicolumn{2}{@{}l}{\textit{Maze and spaces}} \\
\greyrule
maze map & \texttt{large}: 12 columns $\times$ 9 rows of 1\,m cells, walls on the border,
  46 open cells \\
\greyrule
world extent & $x \in [-6, 6]$, $y \in [-4.5, 4.5]$; the origin is the maze centre and row $0$ is
  the top row \\
\greyrule
agent body & 2-D ball, force-actuated, radius $0.1$\,m \\
\greyrule
observation & $4$-d: $x$, $y$, $v_x$, $v_y$. The goal is \emph{not} observed \\
\greyrule
action & $2$-d force, each component clipped to $[-1, 1]$ \\
\blackrule
\multicolumn{2}{@{}l}{\textit{Start and goal}} \\
\greyrule
start cell & row $7$, column $1$ (row $0$ is the top); world position $(-4.5, -3.0)$ \\
\greyrule
goal cell & row $1$, column $10$; world position $(4.5, 3.0)$ \\
\greyrule
reset position noise & $0$. Every episode starts exactly at the start cell's centre and the goal
  sits exactly at its own centre; the only randomness left in an episode is the policy's own
  sampling \\
\greyrule
initial velocity & $(0, 0)$ at every reset \\
\greyrule
goal radius & $0.45$\,m: a step whose distance to the goal is at most this earns reward $1$ \\
\blackrule
\multicolumn{2}{@{}l}{\textit{Episode and reward}} \\
\greyrule
episode length cap & $400$ steps; one step is one $0.01$\,s integrator step \\
\greyrule
task type & continuing: an episode is never terminated at the goal, only truncated at the cap, so
  a copy that reaches the goal keeps collecting reward until the cap \\
\greyrule
reward & sparse: $1$ on a step inside the goal radius, $0$ otherwise; no shift is added \\
\greyrule
discount, extrinsic & $0.999$ \\
\greyrule
discount, intrinsic & $0.99$ (used only by an arm that has an intrinsic bonus) \\
\blackrule
\end{tabular}
\caption{Environment specification for all runs in this section. The discounts are trainer defaults.}
\label{tab:pm-jax-env}
\end{table}

\begin{figure}[H]
\centering
\includegraphics[width=\linewidth]{../../07_reconstruction/development_document/code/2026-06-23-pointmaze-discretation/pointmaze_discretization.pdf}
\caption{Large PointMaze and visit-count grids, reused from the reconstruction development
document. Left: grey cells are walls; the blue circle marks the fixed start at $(-4.5,-3.0)$
and the green star marks the goal at $(4.5,3.0)$. Each position cell is $1$\,m square.
Right: the $10\times10$ velocity bins used with position cells by the oracle bonuses.}
\label{fig:pm-jax-environment}
\end{figure}

\subsection{Train run 1.1: the first algorithm batch}

These historical results were produced before the intrinsic reset correction. At episode
resets, the old implementation continued intrinsic advantages while predicting the next
intrinsic value from the final state. The historical no-bonus control also used weight one
on its intrinsic advantage despite zero intrinsic rewards. The corrected results appear in
Section~\ref{sec:pm-jax-corrected}.
\label{sec:pm-jax-run11}

\begin{table}[H]
\centering
\footnotesize
\setlength{\tabcolsep}{3pt}
\renewcommand{\arraystretch}{1.2}
\resizebox{\textwidth}{!}{%
% Widths: the arm columns are sized by their content, and every spanning cell's width is the sum
% of the columns it covers (11.3 = 3.1+2.3+2.95+2.95, 5.9 = 2.95+2.95).
\begin{tabular}{@{}>{\raggedright\arraybackslash}p{2.9cm}
?>{\raggedright\arraybackslash}p{3.1cm}
?>{\raggedright\arraybackslash}p{2.3cm}
?>{\raggedright\arraybackslash}p{2.95cm}
?>{\raggedright\arraybackslash}p{2.95cm}@{}}
\blackrule
\textbf{Hyperparameter} & \textbf{\texttt{rnd\allowbreak\_next\allowbreak\_state}} &
\textbf{\texttt{no\allowbreak\_exploration}} &
\textbf{\texttt{gt\allowbreak\_position\allowbreak\_velocity\allowbreak\_sqrt}} &
\textbf{\texttt{gt\allowbreak\_position\allowbreak\_velocity\allowbreak\_linear}} \\
& \small(random network distillation) & \small(no intrinsic bonus) &
\small(oracle counts, $1/\sqrt{n}$) & \small(oracle counts, $1/n$) \\
\blackrule
\multicolumn{5}{@{}l}{\textit{Environment and RL algorithm}} \\
\greyrule
environment specification & \multicolumn{4}{>{\raggedright\arraybackslash}p{11.3cm}}{large PointMaze without reset noise (\Cref{tab:pm-jax-env})} \\
\greyrule
base RL algorithm & proximal policy optimisation, this platform's JAX trainer &
  \multicolumn{3}{c}{same} \\
\greyrule
update style & one gradient step per rollout (\texttt{full\allowbreak\_batch}) &
  \multicolumn{3}{c}{same} \\
\greyrule
rollout length $T$ & $128$ steps & \multicolumn{3}{c}{same} \\
\greyrule
environments per copy $N$ & $4$, so one iteration collects $T \times N = 512$ transitions per copy &
  \multicolumn{3}{c}{same} \\
\greyrule
environment steps per copy & $10$M & \multicolumn{3}{c}{same} \\
\blackrule
\multicolumn{5}{@{}l}{\textit{Intrinsic bonus}} \\
\greyrule
bonus signal & the prediction error of a trained network against a fixed random target & N/A &
  \multicolumn{2}{>{\centering\arraybackslash}p{5.9cm}}{the oracle count of the discrete next
  state} \\
\greyrule
bonus input & the next state $(x, y, v_x, v_y)$ & N/A &
  \multicolumn{2}{>{\centering\arraybackslash}p{5.9cm}}{the next state's maze cell and its two
  velocity bins} \\
\greyrule
bonus value & the squared prediction error, whitened by its running standard deviation & N/A &
  $\min(1, n^{-1/2})$ at count~$n$ & $\min(1, n^{-1})$ at count~$n$ \\
\greyrule
what the bonus holds & a predictor network trained with the agent: hidden widths $256$ and $128$;
  $128$ outputs & nothing & \multicolumn{2}{>{\centering\arraybackslash}p{5.9cm}}{one integer count table
  per copy, $9 \times 12 \times 100 = 10{,}800$ entries; a wall or unvisited cell scores $1$} \\
\greyrule
when the bonus updates & every gradient step, as a term of the joint loss & N/A &
  \multicolumn{2}{>{\centering\arraybackslash}p{5.9cm}}{once per rollout: one scatter that counts
  it, one gather that scores it} \\
\greyrule
where the bonus enters & the advantage: $\beta \times$ intrinsic $+\ 2 \times$ extrinsic & N/A &
  \multicolumn{2}{>{\centering\arraybackslash}p{5.9cm}}{the advantage: $\beta \times$ intrinsic
  $+\ 2 \times$ extrinsic} \\
\blackrule
\multicolumn{5}{@{}l}{\textit{Swept knobs}} \\
\greyrule
\textbf{learning rate} & \multicolumn{4}{c}{\bfseries\boldmath $10^{-3}$, $10^{-4}$, $10^{-5}$ --- 3 values} \\
\greyrule
{\bfseries\boldmath intrinsic weight $\beta$} & {\bfseries\boldmath $10^{-5}$ to $10^{5}$ in factor-of-ten steps --- 11 values} & \textbf{N/A} &
  \multicolumn{2}{>{\centering\arraybackslash}p{5.9cm}}{\bfseries\boldmath $10^{-5}$ to $10^{5}$ in factor-of-ten
  steps --- 11 values} \\
\blackrule
\multicolumn{5}{@{}l}{\textit{Run bookkeeping (sweep config counts)}} \\
\greyrule
configs in this column & $3$(learning rate) $\times$ $11$(intrinsic weight) $=33$ &
  $3$(learning rate) $=3$ & $3$(learning rate) $\times$ $11$(intrinsic weight) $=33$ &
  $3$(learning rate) $\times$ $11$(intrinsic weight) $=33$ \\
\greyrule
total configs in this table & \multicolumn{4}{c}{$33+3+33+33=\mathbf{102}$} \\
\greyrule
copies in this column & $33 \times 256 = 8{,}448$ & $3 \times 256 = 768$ &
  $33 \times 256 = 8{,}448$ & $33 \times 256 = 8{,}448$ \\
\greyrule
seeds & \multicolumn{4}{>{\raggedright\arraybackslash}p{11.3cm}}{$256$ copies per configuration,
  with matching seeds across configurations} \\
\blackrule
\multicolumn{5}{@{}l}{\textit{Extension}} \\
\greyrule
first extension & \multicolumn{4}{>{\raggedright\arraybackslash}p{11.3cm}}{each arm's
  best configuration at $10$M steps re-run for $1000$M steps with $1{,}024$ copies, to test
  longer training} \\
\greyrule
second extension & \multicolumn{4}{>{\raggedright\arraybackslash}p{11.3cm}}{the full grids
  for random network distillation and the oracle $1/\sqrt{n}$ bonus re-run for $1000$M steps
  with $128$ copies per configuration, to test whether the best settings change with the budget} \\
\blackrule
\end{tabular}%
}
\caption{Parameters swept in run 1.1 and its two extensions.
``same'' denotes a shared value; ``N/A'' denotes a parameter that does not apply.}
\label{tab:pm-jax-run11-swept}
\end{table}

\paragraph{Run completion.}
All three runs are complete. The original sweep evaluated $102$ configurations at $10$M steps
with $256$ copies each. The first extension, on 2026-08-16--17, tested each arm's selected
configuration at $1000$M steps with $1{,}024$ copies. The second extension, on 2026-08-17--18,
tested all $66$ configurations of random network distillation and the oracle $1/\sqrt{n}$ bonus
at $1000$M steps with $128$ copies each.

\paragraph{Metrics.}
A configuration is one algorithm with a specified learning rate and intrinsic weight.
Rewards are recorded in windows of $200$ training iterations, equivalent to $256$ complete
$400$-step episodes per copy. These windows cover the episode phases evenly.
The tables report means and standard errors over completed copies; the column $N$ gives
that copy count.
\begin{itemize}[leftmargin=*]
\item \textbf{Whole-run reward:} mean episode return over the entire run; used to rank configurations.
\item \textbf{Max reward:} highest window mean, computed by averaging copies first and then
  taking the maximum. Its standard error is measured at that window.
\item \textbf{Last-window reward:} mean episode return in the final recording window.
\item \textbf{Success rate:} fraction of copies that reached the goal at least once.
\item \textbf{Maze-cell coverage:} percentage of the $46$ open cells visited during training.
\item \textbf{Steps to goal:} not recorded.
\end{itemize}

\begin{table}[H]
\centering
\footnotesize
\setlength{\tabcolsep}{3pt}
\renewcommand{\arraystretch}{1.15}
\resizebox{\textwidth}{!}{%
% >>> AUTO-GENERATED TABLE START: pm-jax-run11-results
\begin{tabular}{@{}>{\raggedright\arraybackslash}p{4.4cm} r r r r r r r@{}}
\toprule
\textbf{algorithm --- best configuration} &
\textbf{\shortstack[c]{whole-run\\reward $\uparrow$}} &
\textbf{\shortstack[c]{max\\reward}} &
\textbf{\shortstack[c]{last-window\\reward}} &
\textbf{\shortstack[c]{success\\rate}} &
\textbf{\shortstack[c]{maze-cell\\coverage \%}} &
\textbf{\shortstack[c]{steps\\to goal}} &
\textbf{$N$} \\
\midrule
\multicolumn{8}{@{}l}{\textbf{\texttt{pointmaze\_\allowbreak large\_\allowbreak cont400\_\allowbreak nonoise@1}}} \\
\texttt{gt\_\allowbreak position\_\allowbreak velocity\_\allowbreak sqrt}\newline \small learning rate $10^{-3}$, $\beta = 10^{0}$ & {\boldmath $14.608 \pm 0.981$} & {\boldmath $31.417 \pm 1.916$} & {\boldmath $31.339 \pm 1.915$} & {\boldmath $0.945$} & \underline{$96.51 \pm 0.48$} & not recorded & $256$ \\
\texttt{rnd\_\allowbreak next\_\allowbreak state}\newline \small learning rate $10^{-3}$, $\beta = 10^{1}$ & \underline{$7.866 \pm 0.723$} & \underline{$17.070 \pm 1.464$} & \underline{$17.070 \pm 1.464$} & \underline{$0.926$} & {\boldmath $96.86 \pm 0.48$} & not recorded & $256$ \\
\texttt{gt\_\allowbreak position\_\allowbreak velocity\_\allowbreak linear}\newline \small learning rate $10^{-3}$, $\beta = 10^{0}$ & $2.423 \pm 0.464$ & $5.667 \pm 1.031$ & $5.654 \pm 1.026$ & $0.500$ & $80.06 \pm 1.35$ & not recorded & $256$ \\
\texttt{no\_\allowbreak exploration}\newline \small learning rate $10^{-3}$, no intrinsic weight & $0.005 \pm 0.005$ & $0.026 \pm 0.026$ & $0.019 \pm 0.019$ & $0.004$ & $24.07 \pm 0.89$ & not recorded & $256$ \\
\midrule\midrule
\texttt{rnd\_\allowbreak next\_\allowbreak state}\newline \small learning rate $10^{-3}$, $\beta = 10^{1}$ & {\boldmath $1.795 \pm 0.101$} & \underline{$25.776 \pm 0.955$} & {\boldmath $1.242 \pm 0.286$} & {\boldmath $0.999$} & {\boldmath $99.58 \pm 0.06$} & not recorded & $1{,}024$ \\
\texttt{gt\_\allowbreak position\_\allowbreak velocity\_\allowbreak sqrt}\newline \small learning rate $10^{-3}$, $\beta = 10^{0}$ & \underline{$0.658 \pm 0.024$} & {\boldmath $31.447 \pm 1.053$} & \underline{$0.185 \pm 0.118$} & {\boldmath $0.999$} & \underline{$99.07 \pm 0.06$} & not recorded & $1{,}024$ \\
\texttt{gt\_\allowbreak position\_\allowbreak velocity\_\allowbreak linear}\newline \small learning rate $10^{-3}$, $\beta = 10^{0}$ & $0.331 \pm 0.030$ & $7.181 \pm 0.587$ & $0.171 \pm 0.095$ & \underline{$0.804$} & $94.50 \pm 0.35$ & not recorded & $1{,}024$ \\
\texttt{no\_\allowbreak exploration}\newline \small learning rate $10^{-3}$, no intrinsic weight & $0.001 \pm 0.001$ & $0.088 \pm 0.088$ & $0.000 \pm 0.000$ & $0.016$ & $33.36 \pm 0.53$ & not recorded & $1{,}024$ \\
\midrule\midrule
\texttt{rnd\_\allowbreak next\_\allowbreak state}\newline \small learning rate $10^{-4}$, $\beta = 10^{2}$ & {\boldmath $52.500 \pm 3.368$} & {\boldmath $78.532 \pm 4.113$} & {\boldmath $50.754 \pm 4.851$} & {\boldmath $1.000$} & {\boldmath $99.85 \pm 0.06$} & not recorded & $128$ \\
\texttt{gt\_\allowbreak position\_\allowbreak velocity\_\allowbreak sqrt}\newline \small learning rate $10^{-4}$, $\beta = 10^{1}$ & \underline{$38.199 \pm 3.216$} & \underline{$60.835 \pm 4.816$} & \underline{$33.854 \pm 4.629$} & {\boldmath $1.000$} & {\boldmath $99.85 \pm 0.07$} & not recorded & $128$ \\
\bottomrule
\end{tabular}
% <<< AUTO-GENERATED TABLE END: pm-jax-run11-results
}
\caption{Best configuration per algorithm, ranked by whole-run reward.
The three blocks show the original $10$M-step sweep ($256$ copies), its selected configurations
re-run at $1000$M steps ($1{,}024$ copies), and the full two-arm sweep at $1000$M steps
($128$ copies). Values are means $\pm$ standard errors where shown; $N$ is the number of
completed copies. Best values are \textbf{bold} and second-best values \underline{underlined}
within each block. Higher is better for every recorded metric.}
\label{tab:pm-jax-run11-results}
\end{table}

\begin{figure}[H]
\centering
\includegraphics[width=\textwidth]{\detokenize{../analysis/2026-10-19-12-50_extrinsic-and-raw-intrinsic-reward-curves/plots/historical_extrinsic_selected.pdf}}
\caption{Mean extrinsic reward per episode over $1000$M steps per copy, with bands of
one standard error. Dashed curves use configurations selected at $10$M steps ($1{,}024$ copies);
solid curves use configurations selected at $1000$M steps ($128$ copies). Each point averages
$256$ episodes per copy. The step axis is logarithmic. The legend gives the learning rate and
intrinsic weight.}
\label{fig:pm-jax-run11-curves}
\end{figure}

\paragraph{Findings.}
\begin{enumerate}[label=\textbf{F\arabic*},leftmargin=*,labelsep=0.6em]
\item \textbf{Distillation improves reward at the two larger learning rates.} Random network distillation
  improves reward over training without a bonus at $10^{-3}$ and $10^{-4}$. At $10^{-5}$,
  both have zero reward, although distillation achieves higher coverage.
\item \textbf{Oracle bonuses do not consistently outperform distillation.} At $10$M steps, the oracle $1/\sqrt{n}$ bonus has higher reward
  than distillation, but $1/n$ has lower reward. The winning oracle's coverage is not clearly
  different from distillation's.
\item \textbf{The slower-decaying oracle bonus performs better on both metrics.} At $10$M steps, $1/\sqrt{n}$ achieves both higher coverage and
  higher reward than $1/n$.
\item \textbf{Reward peaks at an intermediate intrinsic weight.} At learning rate $10^{-3}$, reward has an interior peak
  and declines at larger intrinsic weights. Very small weights yield almost no reward,
  rather than a clear increasing sequence.
\item \textbf{The best settings depend on the training budget.} The $10$M-step sweep selects
  learning rate $10^{-3}$; the $1000$M-step sweep selects $10^{-4}$ and an intrinsic weight
  ten times larger for both re-swept arms. At the long budget, these settings achieve about
  $29$ times the reward for distillation and $58$ times the reward for the oracle bonus
  compared with the settings selected at $10$M steps.
\item \textbf{Distillation has higher reward at the long budget.} With settings selected at
  $1000$M steps, its whole-run reward is $52.500 \pm 3.368$, compared with
  $38.199 \pm 3.216$ for the oracle $1/\sqrt{n}$ bonus. Both reach the goal in all $128$
  copies and achieve $99.85\%$ coverage.
\end{enumerate}

\paragraph{Curve shapes.}
Not every configuration rises and then falls. Of the $66$ configurations in the second
extension, $20$ rise then fall, $16$ never reach the goal, $14$ are still rising, $13$ decline
from the start, and $3$ rise then flatten. Labels remain unchanged under all nine classification
threshold variants for $53$ configurations, including the ten with highest whole-run reward.
\Cref{tab:pm-jax-run11-shapes} shows the best non-falling configuration of each arm and its
overall best for comparison.

\Cref{fig:pm-jax-run11-shapes} compares learning rates. At $10^{-3}$, the selected curves peak within
$15$M steps and then decline; at $10^{-4}$, they peak near $200$M steps at higher returns.
Both arms show this pattern.

\begin{table}[H]
\centering
\footnotesize
\setlength{\tabcolsep}{3.5pt}
\renewcommand{\arraystretch}{1.15}
\begin{tabular}{@{}>{\raggedright\arraybackslash}p{3.4cm} r r >{\raggedright\arraybackslash}p{1.9cm} r r r@{}}
\toprule
\textbf{algorithm} & \textbf{\shortstack[c]{learning\\rate}} &
\textbf{\shortstack[c]{intrinsic\\weight}} & \textbf{aggregate shape} &
\textbf{\shortstack[c]{whole-run\\reward $\uparrow$}} & \textbf{\shortstack[c]{rank of\\$66$}} &
\textbf{\shortstack[c]{copies of $128$\\sharing it}} \\
\midrule
\texttt{rnd\_\allowbreak next\_\allowbreak state} & $10^{-5}$ & $10^{1}$ & still rising &
  $27.778$ & $6$ & $55$ \\
\texttt{gt\_\allowbreak position\_\allowbreak velocity\_\allowbreak sqrt} & $10^{-4}$ & $10^{2}$ &
  rise then plateau & $10.391$ & $11$ & $6$ \\
\midrule
\multicolumn{7}{@{}l}{\textit{for comparison, the same arms' best configurations, which do rise then
  fall}} \\
\texttt{rnd\_\allowbreak next\_\allowbreak state} & $10^{-4}$ & $10^{2}$ & rise then fall &
  $52.500$ & $1$ & $76$ \\
\texttt{gt\_\allowbreak position\_\allowbreak velocity\_\allowbreak sqrt} & $10^{-4}$ & $10^{1}$ &
  rise then fall & $38.199$ & $3$ & $84$ \\
\bottomrule
\end{tabular}
\caption{Best configuration of each arm whose mean curve does not rise
then fall, followed by the overall best configurations. Rank is by whole-run reward across
all $66$ configurations. The last column counts copies sharing the mean curve's shape:
individual copies can behave differently. The oracle plateau label changes to ``still rising''
under two of nine threshold variants; $90$ of its copies individually rise then fall.}
\label{tab:pm-jax-run11-shapes}
\end{table}

\paragraph{Reward quantities.}
The existing training curves show extrinsic reward only. Let $R_{\mathrm{ext}}$ be the
extrinsic reward summed over a 400-step episode, and let $R_{\mathrm{int,raw}}$ be the raw
intrinsic reward summed over the same episode. The additional plots show their unweighted
sum, $R_{\mathrm{sum}}=R_{\mathrm{ext}}+R_{\mathrm{int,raw}}$.
The logged intrinsic mean is per step, so it is multiplied by 400 before addition.
These sums use the raw bonus before normalization and do not apply the intrinsic weight.
The trainer combines separately discounted task and normalized intrinsic advantages; its
normalization values were not saved. Its weighted, normalized training total therefore
cannot be reconstructed from these records.

\begingroup
\renewcommand{\thefigure}{\thesection.3.1}
\renewcommand{\theHfigure}{\thesection.3.1}
\begin{figure}[H]
\centering
\includegraphics[width=\textwidth]{\detokenize{../analysis/2026-10-19-12-50_extrinsic-and-raw-intrinsic-reward-curves/plots/historical_extrinsic.pdf}}
\caption{Extrinsic reward per episode for each arm at learning rates $10^{-3}$ and $10^{-4}$, using the intrinsic weight
with the highest peak in the copy-averaged curve. Orange curves add random network distillation
at fixed intrinsic weight $\beta=10^0=1$ for both learning rates. Markers show the peaks of the
selected configurations; bands show one standard error over $128$ copies. Dotted and solid lines
identify learning rates $10^{-3}$ and $10^{-4}$, respectively. Selection by peak reward can differ
from selection by whole-run reward in \Cref{tab:pm-jax-run11-results}.}
\label{fig:pm-jax-run11-shapes}
\end{figure}
\endgroup

\begingroup
\addtocounter{figure}{-1}
\renewcommand{\thefigure}{\thesection.3.2}
\renewcommand{\theHfigure}{\thesection.3.2}
\begin{figure}[H]
\centering
\includegraphics[width=\textwidth]{\detokenize{../analysis/2026-10-19-12-50_extrinsic-and-raw-intrinsic-reward-curves/plots/historical_total.pdf}}
\caption{Total recorded reward per episode: extrinsic plus raw intrinsic, unweighted,
for the same historical settings as Figure~\ref{fig:pm-jax-run11-shapes}. Colors and line
styles match that figure. Each raw intrinsic mean per step is multiplied by 400 before
addition to extrinsic reward per episode. Bands show one standard error of the combined
reward across 128 copies, including covariance between its two components. This raw sum
does not reconstruct the weighted, normalized quantity used in policy updates. The vertical
axis is logarithmic above 1 and linear between $-1$ and $1$.}
\label{fig:pm-jax-run11-total}
\end{figure}
\endgroup



\subsection{Train run 1.2: corrected intrinsic returns}
\label{sec:pm-jax-corrected}
At a 400-step reset, intrinsic advantages now continue into the next episode using the
intrinsic value of the reset observation. The immediate intrinsic bonus still scores the final
transition observation. Task-return bootstrapping and task advantage boundaries retain their
previous definitions. This matches CleanRL's continuing intrinsic-return convention.
For the no-bonus control, the intrinsic advantage now has weight zero, so its policy update
uses only task advantages. The shared two-head critic architecture is retained.

The 11 unique settings shown in Figures~\ref{fig:pm-jax-run11-curves} and
\ref{fig:pm-jax-run11-shapes} are retrained for 100M steps with 128 paired seeds per setting.
The learning rate is denoted by $\alpha$, and the weight on intrinsic advantage by $\beta$.
For the oracle bonuses, $n$ denotes the visit count of the discretized position and velocity state.
Only the listed pairs are trained; the experiment does not take their Cartesian product.

\begin{table}[H]
\centering
\small
\setlength{\tabcolsep}{3pt}
\renewcommand{\arraystretch}{1.14}
\begin{tabular}{@{}p{2.8cm}?p{2.9cm}?p{2.6cm}?p{2.6cm}?p{2.1cm}@{}}
\blackrule
\textbf{hyperparameter} & \textbf{random network distillation} &
\textbf{oracle counts $1/\sqrt{n}$} & \textbf{oracle counts $1/n$} & \textbf{no bonus} \\
\blackrule
\multicolumn{5}{@{}l}{\textit{Environment and algorithm}} \\
\greyrule
environment & \multicolumn{4}{c}{PointMaze large, no reset noise (Table~\ref{tab:pm-jax-env})} \\
\greyrule
training budget per copy & \multicolumn{4}{c}{100M environment steps} \\
\greyrule
update schedule & \multicolumn{4}{c}{one full-batch gradient step per rollout} \\
\greyrule
rollout and environments & \multicolumn{4}{c}{128 steps, four environments per copy} \\
\greyrule
seeds per setting & \multicolumn{4}{c}{128; the same seed indices across all settings} \\
\greyrule
other settings & \multicolumn{4}{c}{as in Table~\ref{tab:pm-jax-run11-swept}; corrected intrinsic reset value} \\
\blackrule
\multicolumn{5}{@{}l}{\textit{Selected parameters}} \\
\greyrule
{\bfseries\boldmath learning rate $\alpha$} & {\bfseries\boldmath $10^{-3},10^{-4}$} &
{\bfseries\boldmath $10^{-3},10^{-4}$} & {\bfseries\boldmath $10^{-3}$} & {\bfseries\boldmath $10^{-3}$} \\
\greyrule
{\bfseries\boldmath intrinsic weight $\beta$} & {\bfseries\boldmath $1,10,100,1000$} &
{\bfseries\boldmath $1,10$} & {\bfseries\boldmath $1$} & \textbf{N/A} \\
\greyrule
{\bfseries\boldmath trained $(\alpha,\beta)$ pairs} &
{\bfseries\boldmath $(10^{-3},1)$\newline $(10^{-3},10)$\newline $(10^{-3},100)$\newline
$(10^{-4},1)$\newline $(10^{-4},100)$\newline $(10^{-4},1000)$} &
{\bfseries\boldmath $(10^{-3},1)$\newline $(10^{-4},1)$\newline $(10^{-4},10)$} &
{\bfseries\boldmath $(10^{-3},1)$} & {\bfseries\boldmath $10^{-3}$, no weight} \\
\blackrule
\multicolumn{5}{@{}l}{\textit{Configuration counts}} \\
\greyrule
configurations & 6 selected pairs & 3 selected pairs & 1 selected pair & 1 learning rate \\
\greyrule
total configurations & \multicolumn{4}{c}{$6+3+1+1=11$} \\
\blackrule
\end{tabular}
\caption{The settings retrained after correcting intrinsic returns. Each listed setting uses 128 independent copies.}
\label{tab:pm-jax-corrected-swept}
\end{table}

\paragraph{Metrics.}
Whole-run reward is each copy's mean task return over complete episode windows, then averaged
over copies. Final-window reward uses the last 160 episodes per copy. Goal-reaching is the
fraction of copies that received task reward at least once during training. Coverage is the
fraction of the 46 open maze cells visited. Reported uncertainties are standard errors over
128 copies. The settings are fixed from the historical figures, rather than selected again
from a full grid. A 100M-step run cannot establish the later peaks seen at 1000M steps.
The number of completed copies is denoted by $N$.

\begin{table}[H]
\centering
\small
\setlength{\tabcolsep}{3pt}
% Generator: ../scripts/write_corrected_report.py; reads the run's summary.json.
% >>> AUTO-GENERATED TABLE START: pm-jax-corrected-results
\begin{tabular}{@{}p{3.4cm}rrrrr@{}}
\toprule
\textbf{algorithm and settings} & \shortstack{whole-run\\reward $\uparrow$} &\shortstack{final-window\\reward} & \shortstack{ever\\reached goal} &\shortstack{coverage\\(\%)} & $N$ \\
\midrule
random network distillation\newline $\alpha=10^{-3}$, $\beta=10^{2}$ & {\boldmath $24.163 \pm 1.753$} & {\boldmath $33.045 \pm 3.916$} & {\boldmath $1.000$} & {\boldmath $99.304 \pm 0.288$} & 128 \\
oracle counts $1/\sqrt{n}$\newline $\alpha=10^{-3}$, $\beta=10^{0}$ & \underline{$17.875 \pm 1.676$} & \underline{$18.295 \pm 3.374$} & {\boldmath $1.000$} & \underline{$98.998 \pm 0.174$} & 128 \\
oracle counts $1/n$\newline $\alpha=10^{-3}$, $\beta=10^{0}$ & $6.849 \pm 1.398$ & $11.882 \pm 2.932$ & \underline{$0.734$} & $91.372 \pm 1.218$ & 128 \\
no bonus\newline $\alpha=10^{-3}$ & $0.000 \pm 0.000$ & $0.000 \pm 0.000$ & $0.008$ & $26.919 \pm 1.213$ & 128 \\
\bottomrule
\end{tabular}
% <<< AUTO-GENERATED TABLE END: pm-jax-corrected-results
\caption{Corrected 100M-step results at each algorithm's highest whole-run reward among the
retrained settings. Values are means $\pm$ standard errors where shown. Selection is among the
11 settings in Table~\ref{tab:pm-jax-corrected-swept}, rather than the full historical grid.}
\label{tab:pm-jax-corrected-results}
\end{table}

\paragraph{Findings.}
At intrinsic weight $1000$, random network distillation reaches $99.88\%$ coverage with
whole-run reward $0.44$. Without a bonus, only one of the $128$ copies ever reaches the goal;
its whole-run mean reward is $3.44\times10^{-7}$, which rounds to zero in the table.


\begingroup
\renewcommand{\thefigure}{\thesection.4.1}
\renewcommand{\theHfigure}{\thesection.4.1}
\begin{figure}[H]
\centering
\includegraphics[width=\textwidth]{\detokenize{../analysis/2026-10-19-12-50_extrinsic-and-raw-intrinsic-reward-curves/plots/corrected_extrinsic.pdf}}
\caption{Extrinsic reward per episode during corrected 100M-step training for the six settings in historical
Figure~\ref{fig:pm-jax-run11-shapes}. Colors retain their original meaning: random network
distillation is black, oracle $1/\sqrt{n}$ counts are blue, and random network distillation at
fixed intrinsic weight $\beta=1$ is orange. Dotted and solid lines identify learning rates
$10^{-3}$ and $10^{-4}$, respectively. Bands show one standard error over 128 copies.
The settings are held fixed from the historical figure; no new peak winners are selected.}
\label{fig:pm-jax-corrected-curves}
\end{figure}
\endgroup

\begingroup
\addtocounter{figure}{-1}
\renewcommand{\thefigure}{\thesection.4.2}
\renewcommand{\theHfigure}{\thesection.4.2}
\begin{figure}[H]
\centering
\includegraphics[width=\textwidth]{\detokenize{../analysis/2026-10-19-12-50_extrinsic-and-raw-intrinsic-reward-curves/plots/corrected_total.pdf}}
\caption{Total recorded reward per episode: extrinsic plus raw intrinsic, unweighted,
during corrected 100M-step training for the same six settings as
Figure~\ref{fig:pm-jax-corrected-curves}. Colors and line styles match that figure.
The reward definition and limitation are the same as in Figure~\ref{fig:pm-jax-run11-total}:
raw intrinsic reward is added before normalization and weighting. Bands show one standard
error of each copy's combined reward over 128 copies. These plots describe the recorded
reward trajectories; they do not establish why extrinsic reward rises and then falls. The
vertical axis is logarithmic above 1 and linear between $-1$ and $1$.}
\label{fig:pm-jax-corrected-total}
\end{figure}
\endgroup


\subsection{Training speed on jaguar03}
\studyunreadbegin % REPORT-STATUS
\label{sec:pm-jax-speed-jaguar03}
The reward-decline study below runs many short units on jaguar03's eight RTX A4500 cards. Ten
execution ideas were timed there on 2026-10-19 against one baseline: a random network
distillation program of $128$ copies (learning rate $10^{-3}$, intrinsic weight $100$, full-batch
updates, coverage tracking on). Each idea ran twice, on two different cards, while a baseline ran
on a third card at the same time. Let $r$ be the steps per second one copy gets; a unit of
$S$ million steps per copy then takes $S \times 10^{6}/(3600\,r)$ hours. ``Total against
baseline'' divides an idea's total steps per second by the baseline's.

\begin{table}[H]
\studyfloatcolor % REPORT-STATUS
\centering
\footnotesize
\setlength{\tabcolsep}{3pt}
\renewcommand{\arraystretch}{1.15}
\resizebox{\textwidth}{!}{%
% Generator: ../scripts/write_decline_study_tables.py; reads the benchmark's analysis_rows.json.
% >>> AUTO-GENERATED TABLE START: pm-jax-speed-jaguar03
\begin{tabular}{@{}>{\raggedright\arraybackslash}p{4.3cm} r r r r r r r >{\raggedright\arraybackslash}p{2.6cm}@{}}
\toprule
\textbf{idea} & \textbf{\shortstack[c]{copies}} & \textbf{\shortstack[c]{seconds per\\iteration}} & \textbf{\shortstack[c]{total steps\\per second (M)}} & \textbf{\shortstack[c]{steps per second\\per copy (k)}} & \textbf{\shortstack[c]{hours per M\\steps per copy}} & \textbf{\shortstack[c]{total against\\baseline $\uparrow$}} & \textbf{\shortstack[c]{setup\\seconds}} & \textbf{\shortstack[c]{copies identical\\to baseline}} \\
\midrule
02 eight settings in one program & 1{,}024 & 0.0505 & \textbf{10.39} & 10.1 & 0.0274 & \textbf{1.44} & 84 & 0 of 128 \\
01 four settings in one program & 512 & 0.0269 & \underline{9.75} & 19.0 & 0.0146 & \underline{1.35} & 68 & 0 of 128 \\
07 XLA latency-hiding scheduler on & 128 & 0.0090 & 7.25 & \textbf{56.6} & \textbf{0.0049} & 1.00 & 59 & 0 of 128 \\
08 compilation cache, first run & 128 & 0.0091 & 7.24 & \textbf{56.6} & \textbf{0.0049} & 1.00 & 57 & 5 of 128 \\
baseline: one setting, one program per card & 128 & 0.0091 & 7.22 & \underline{56.4} & \underline{0.0049} & 1.00 & 55 & 3 of 128 \\
09 two processors instead of eight & 128 & 0.0091 & 7.22 & \underline{56.4} & \underline{0.0049} & 1.00 & 79 & 91 of 128 \\
08 compilation cache, second run & 128 & 0.0091 & 7.22 & \underline{56.4} & \underline{0.0049} & 1.00 & \textbf{7} & 5 of 128 \\
04 four programs share a card & 512 & 0.0381 & 6.88 & 13.4 & 0.0207 & 0.95 & 85 & 1 of 128 \\
03 two programs share a card & 256 & 0.0192 & 6.83 & 26.7 & 0.0104 & 0.95 & 66 & 91 of 128 \\
06 XLA command buffers also for loops & 128 & 0.0101 & 6.49 & 50.7 & 0.0055 & 0.90 & 58 & 26 of 128 \\
10 the study's diagnostics on & 128 & 0.0105 & 6.23 & 48.7 & 0.0057 & 0.86 & \underline{50} & 0 of 128 \\
05 XLA command buffers off & 128 & 0.0128 & 5.17 & 40.4 & 0.0069 & 0.72 & 60 & 91 of 128 \\
\bottomrule
\end{tabular}
% <<< AUTO-GENERATED TABLE END: pm-jax-speed-jaguar03
}
\caption{Speed of the ten execution ideas on one RTX A4500, sorted by total steps per second
against the baseline. M is millions and k thousands of environment steps. Setup is building,
priming and compiling before the first timed iteration. The last column counts, of setting~0's
$128$ copies, those whose parameters after ten iterations match a baseline run bit for bit; the
baseline row compares two baseline runs on two cards. Best values are \textbf{bold} and second
best \underline{underlined}; higher is better for rates and lower for hours and setup.}
\label{tab:pm-jax-speed-jaguar03}
\end{table}

\begin{figure}[H]
\studyfloatcolor % REPORT-STATUS
\centering
\includegraphics[width=\textwidth]{\detokenize{../benchmark_runs/2026-10-19-15-46_jaguar03_rtx-a4500_random-network-distillation_copies-128-per-setting_rollout-128_environments-4_full-batch_speed-10-ideas-round-2/plots/speed_ideas.pdf}}
\caption{Two ways of putting more copies on one RTX A4500: settings packed into one program
(each setting $128$ copies), and several $128$-copy programs sharing the card. Left: total steps
per second; right: steps per second per copy.}
\label{fig:pm-jax-speed-jaguar03}
\end{figure}

\paragraph{Findings.}
\begin{enumerate}[label=\textbf{S\arabic*},leftmargin=*,labelsep=0.6em]
\item \textbf{Packing settings is the only idea that raises a card's throughput.} Four settings
  in one program give $1.35$ times the baseline's total, eight give $1.44$ times; each setting
  then advances $3.0$ and $5.6$ times more slowly.
\item \textbf{Sharing a card between programs lowers the total} to $0.95$ times, for two and for
  four programs.
\item \textbf{XLA's default command buffers are already the fastest choice.} Turning them off
  gives $0.72$ times; extending them to loops gives $0.90$ times. The latency-hiding scheduler
  changes nothing.
\item \textbf{A compilation cache cuts setup from 57 to 7 seconds} on the second run of a
  program; two processors instead of eight leave the rate unchanged but slow setup to 79 seconds.
\item \textbf{The study's diagnostics cost 14\%} of the rate.
\item \textbf{Runs on these cards are reproducible only to rounding.} Two runs of the same
  program on two cards keep $3$ of $128$ copies bit-identical after ten iterations; the rest
  differ by at most $2.8\times10^{-4}$. A program with a different copy count draws different
  policy noise for the same copy, because the noise array's shape includes the copy count; so
  copies are paired only between programs of equal copy count.
\end{enumerate}

\paragraph{Choices for the study.} Every unit packs the four intrinsic weights $1$, $10$, $100$
and $1000$ into one program of $4 \times 128$ copies, so every unit of a batch has the same copy
count and copy $k$ meets the same noise and environments in every idea. Jobs share one
compilation cache, ask for eight processors, and keep the diagnostics on.
\studyunreadend % REPORT-STATUS

\subsection{Choosing a small environment for the reward-decline study}
\studyunreadbegin % REPORT-STATUS
\label{sec:pm-jax-screening}
Candidate fixes are tested first on the smallest maze that still shows the problem. Three smaller
PointMaze maps were screened against the large one, with the same ball dynamics, goal radius,
sparse reward and continuing task as \Cref{tab:pm-jax-env}; only the map, the start, the goal and
the episode cap differ (\Cref{tab:pm-jax-screening-env}, \Cref{fig:pm-jax-screening-maps}). The
medium maze is screened at two caps: $400$ steps, and $250$, which keeps the large maze's ratio of
travel to cap ($10$ cell moves against $17$).

\begin{table}[H]
\studyfloatcolor % REPORT-STATUS
\centering
\small
\setlength{\tabcolsep}{3pt}
\renewcommand{\arraystretch}{1.15}
\begin{tabular}{@{}>{\raggedright\arraybackslash}p{2.8cm}?>{\raggedright\arraybackslash}p{2.3cm}?>{\raggedright\arraybackslash}p{2.6cm}?>{\raggedright\arraybackslash}p{2.6cm}?>{\raggedright\arraybackslash}p{2.6cm}@{}}
\blackrule
\textbf{parameter} & \textbf{U-maze} & \textbf{medium, cap 250} & \textbf{medium, cap 400} &
\textbf{large, cap 400} \\
\blackrule
open cells & $7$ & \multicolumn{2}{c?}{$26$} & $46$ \\
\greyrule
start cell (row, column) & $(3,1)$ & \multicolumn{2}{c?}{$(6,1)$} & $(7,1)$ \\
\greyrule
goal cell & $(1,1)$ & \multicolumn{2}{c?}{$(1,6)$} & $(1,10)$ \\
\greyrule
shortest route & $6$ cell moves & \multicolumn{2}{c?}{$10$ cell moves} & $17$ cell moves \\
\greyrule
episode cap & $150$ steps & $250$ steps & $400$ steps & $400$ steps \\
\greyrule
steps per copy & \multicolumn{3}{c?}{$30$M} & $50$M \\
\blackrule
\end{tabular}
\caption{The screened maps. Everything not listed is as in \Cref{tab:pm-jax-env}.}
\label{tab:pm-jax-screening-env}
\end{table}

\begin{figure}[H]
\studyfloatcolor % REPORT-STATUS
\centering
\includegraphics[width=\textwidth]{\detokenize{../runs/2026-10-19-16-18_map-screening_umaze-cap150_medium-cap250_medium-cap400_large-cap400_rnd-weights-1-10-100-1000_visit-count-sqrt-weight-1_no-bonus_learning-rate-1e-3_steps-30M-large-50M_seeds-0-127/analysis/plots/screened_maps.pdf}}
\caption{The screened maps, to scale with $1$\,m cells. Grey cells are walls; the blue disc
marks the start and the green star and circle the goal and its $0.45$\,m radius.}
\label{fig:pm-jax-screening-maps}
\end{figure}

\begin{table}[H]
\studyfloatcolor % REPORT-STATUS
\centering
\small
\setlength{\tabcolsep}{3pt}
\renewcommand{\arraystretch}{1.14}
\begin{tabular}{@{}>{\raggedright\arraybackslash}p{3.1cm}?>{\raggedright\arraybackslash}p{3.7cm}?>{\raggedright\arraybackslash}p{3.2cm}?>{\raggedright\arraybackslash}p{3.2cm}@{}}
\blackrule
\textbf{hyperparameter} & \textbf{random network distillation} &
\textbf{oracle counts $1/\sqrt{n}$} & \textbf{no bonus} \\
\blackrule
\multicolumn{4}{@{}l}{\textit{Environment and algorithm}} \\
\greyrule
{\bfseries\boldmath environment} & \multicolumn{3}{>{\centering\arraybackslash}p{10.3cm}}{\bfseries\boldmath
U-maze cap 150; medium cap 250; medium cap 400; large cap 400 --- 4 values} \\
\greyrule
learning rate & \multicolumn{3}{>{\centering\arraybackslash}p{10.3cm}}{$10^{-3}$, annealed linearly to zero over the run} \\
\greyrule
update and rollout & \multicolumn{3}{>{\centering\arraybackslash}p{10.3cm}}{one full-batch step per rollout of $128$ steps $\times$ $4$
environments} \\
\greyrule
seeds & \multicolumn{3}{c}{$128$ copies per configuration, seeds $0$--$127$} \\
\blackrule
\multicolumn{4}{@{}l}{\textit{Intrinsic bonus}} \\
\greyrule
{\bfseries\boldmath intrinsic weight $\beta$} & {\bfseries\boldmath $1$, $10$, $100$, $1000$ ---
4 values} & {\bfseries\boldmath $1$} & \textbf{N/A} \\
\blackrule
\multicolumn{4}{@{}l}{\textit{Configuration counts}} \\
\greyrule
configurations & $4$(maps) $\times$ $4$(weights) $=16$ & $4$(maps) $=4$ & $4$(maps) $=4$ \\
\greyrule
total configurations & \multicolumn{3}{c}{$16+4+4=24$} \\
\blackrule
\end{tabular}
\caption{The screening's settings. Bold rows are swept.}
\label{tab:pm-jax-screening-swept}
\end{table}

\paragraph{Metrics.}
Episode returns are counted over the episodes that end inside each $250$-iteration recording
window. Let $G_{c,w}$ be the mean extrinsic return of the episodes copy $c$ ended in window $w$,
and $\bar G_w$ its mean over the configuration's copies. Then the peak is $\max_w \bar G_w$; the
final-10\% return is the mean over copies of each copy's return in the last tenth of the windows;
``final / peak'' is their ratio; ``copies ever at the goal'' is the share of copies with at least
one episode inside the goal radius; and ``final share of episodes at goal'' is the share of the
last tenth's episodes that entered the goal radius. Uncertainties are standard errors over
copies.

The rules, fixed before the runs: the chosen map is the smallest on which (1) a policy without a
bonus reaches the goal in at most $25\%$ of copies, so exploration is needed; (2) the random
network distillation setting with the highest whole-run return keeps at most $60\%$ of its peak in
the final tenth, with its peak in the first half of the run; and (3) the oracle $1/\sqrt{n}$ bonus
keeps at most $60\%$ of its peak. If no smaller map passes, the study uses the large maze at
learning rate $10^{-3}$ for $50$M steps.

\paragraph{Run status.} Submitted 2026-10-19 16:22 PT on twelve RTX 2080 Ti cards; the results
tables and curves are added when the last unit ends.
\studyunreadend % REPORT-STATUS

\subsection{Reward-decline batch 1: a control and ten ideas}
\studyunreadbegin % REPORT-STATUS
\label{sec:pm-jax-decline-batch1}
Each idea changes one thing against a shared control. Three change the environment's reward
structure (\Cref{tab:pm-jax-batch1-env-ideas}); seven change the algorithm
(\Cref{tab:pm-jax-batch1-alg-ideas}). Every unit runs the four intrinsic weights $1$, $10$, $100$
and $1000$ as one program of $4\times128$ copies, so copy $k$ meets the same environments and
policy noise in every unit. The learning rate anneals linearly over $1000$M steps, so it stays
above $97.5\%$ of $10^{-3}$ during these $50$M steps, as in the $1000$M-step runs of
\Cref{fig:pm-jax-run11-shapes} where the fall was first seen; annealing over the run's own $50$M
steps would hide part of the fall, as in \Cref{fig:pm-jax-corrected-curves}.

Define $k$ as the iteration index and $N=97{,}657$ as the run's iterations; $\beta$ as the
intrinsic weight; $s_c$ as copy $c$'s running share of ended episodes that reached the goal,
moved each iteration by $0.05$ times the difference to the share among that iteration's ended
episodes; and $T_{\max}=400$ as the episode cap.

\begin{table}[H]
\studyfloatcolor % REPORT-STATUS
\centering
\small
\setlength{\tabcolsep}{3pt}
\renewcommand{\arraystretch}{1.14}
\begin{tabular}{@{}>{\raggedright\arraybackslash}p{2.8cm}?>{\raggedright\arraybackslash}p{2.5cm}?>{\raggedright\arraybackslash}p{2.5cm}?>{\raggedright\arraybackslash}p{2.5cm}?>{\raggedright\arraybackslash}p{2.6cm}@{}}
\blackrule
\textbf{hyperparameter} & \textbf{control} & \textbf{E1} & \textbf{E2} & \textbf{E3} \\
& \small(original task) & \small(end at the goal) & \small(end at the goal, step cost) &
\small(end at the goal, lump sum) \\
\blackrule
\multicolumn{5}{@{}l}{\textit{Environment and algorithm}} \\
\greyrule
environment & \multicolumn{4}{>{\centering\arraybackslash}p{10.5cm}}{large PointMaze,
start $(7,1)$, goal $(1,10)$, cap $400$ steps (\Cref{tab:pm-jax-env})} \\
\greyrule
algorithm & \multicolumn{4}{>{\centering\arraybackslash}p{10.5cm}}{random network distillation
with proximal policy optimisation; one full-batch step per rollout of $128$ steps $\times$ $4$
environments; extrinsic discount $0.999$, intrinsic $0.99$; extrinsic weight $2$} \\
\greyrule
learning rate & \multicolumn{4}{>{\centering\arraybackslash}p{10.5cm}}{$10^{-3}$, annealed
linearly over $1000$M steps} \\
\greyrule
training budget & \multicolumn{4}{c}{$50$M environment steps per copy} \\
\greyrule
seeds & \multicolumn{4}{c}{$128$ copies per configuration, seeds $0$--$127$, paired} \\
\greyrule
{\bfseries\boldmath intrinsic weight $\beta$} &
\multicolumn{4}{>{\centering\arraybackslash}p{10.5cm}}{\bfseries\boldmath $1$, $10$, $100$,
$1000$ in one program --- 4 values} \\
\blackrule
\multicolumn{5}{@{}l}{\textit{Reward structure}} \\
\greyrule
{\bfseries\boldmath episode ends} & {\bfseries\boldmath only at the cap} &
\multicolumn{3}{>{\centering\arraybackslash}p{7.7cm}}{\bfseries\boldmath on the first step inside
the goal radius, or at the cap} \\
\greyrule
{\bfseries\boldmath reward inside the goal radius} & {\bfseries\boldmath $1$ on every such step}
& {\bfseries\boldmath $1$, once} & {\bfseries\boldmath $0$, once} &
{\bfseries\boldmath $T_{\max}-t+1$ at arrival step $t$} \\
\greyrule
{\bfseries\boldmath reward on other steps} & {\bfseries\boldmath $0$} & {\bfseries\boldmath $0$}
& {\bfseries\boldmath $-1$} & {\bfseries\boldmath $0$} \\
\blackrule
\multicolumn{5}{@{}l}{\textit{Configuration counts}} \\
\greyrule
configurations & $4$(weights) & $4$(weights) & $4$(weights) & $4$(weights) \\
\blackrule
\end{tabular}
\caption{Batch 1, the control and the three environment-reward ideas. Bold rows are swept. E3
pays at arrival what staying until the cap would have earned in the control's task.}
\label{tab:pm-jax-batch1-env-ideas}
\end{table}

\begin{table}[H]
\studyfloatcolor % REPORT-STATUS
\centering
\small
\setlength{\tabcolsep}{2.5pt}
\renewcommand{\arraystretch}{1.14}
\resizebox{\textwidth}{!}{%
\begin{tabular}{@{}>{\raggedright\arraybackslash}p{2.7cm}?>{\raggedright\arraybackslash}p{1.9cm}?>{\raggedright\arraybackslash}p{1.9cm}?>{\raggedright\arraybackslash}p{1.9cm}?>{\raggedright\arraybackslash}p{1.6cm}?>{\raggedright\arraybackslash}p{2.0cm}?>{\raggedright\arraybackslash}p{1.6cm}?>{\raggedright\arraybackslash}p{1.9cm}@{}}
\blackrule
\textbf{hyperparameter} & \textbf{A1} & \textbf{A2} & \textbf{A3} & \textbf{A4} & \textbf{A5} &
\textbf{A6} & \textbf{A7} \\
& \small(weight decay) & \small(goal-rate gate) & \small(clipped PPO) & \small(entropy) &
\small(task-reward scaling) & \small(discount) & \small(clip per network) \\
\blackrule
everything else & \multicolumn{7}{c}{as the control in \Cref{tab:pm-jax-batch1-env-ideas},
including $\beta\in\{1,10,100,1000\}$} \\
\blackrule
{\bfseries\boldmath intrinsic weight used} & {\bfseries\boldmath $\beta\,(1-\frac{k-1}{N})$} &
{\bfseries\boldmath $\beta\,(1-s_c)$} & {\bfseries\boldmath $\beta$} & {\bfseries\boldmath $\beta$}
& {\bfseries\boldmath $\beta$} & {\bfseries\boldmath $\beta$} & {\bfseries\boldmath $\beta$} \\
\greyrule
{\bfseries\boldmath update} & \textbf{one step} & \textbf{one step} & {\bfseries\boldmath $4$
epochs $\times$ $4$ minibatches, clip $0.2$} & \textbf{one step} & \textbf{one step} &
\textbf{one step} & \textbf{one step} \\
\greyrule
{\bfseries\boldmath entropy coefficient} & {\bfseries\boldmath $0$} & {\bfseries\boldmath $0$} &
{\bfseries\boldmath $0$} & {\bfseries\boldmath $0.01$} & {\bfseries\boldmath $0$} &
{\bfseries\boldmath $0$} & {\bfseries\boldmath $0$} \\
\greyrule
{\bfseries\boldmath task reward used} & \textbf{as paid} & \textbf{as paid} & \textbf{as paid} &
\textbf{as paid} & {\bfseries\boldmath divided by $\max(1,\sigma)$} & \textbf{as paid} &
\textbf{as paid} \\
\greyrule
{\bfseries\boldmath extrinsic discount} & {\bfseries\boldmath $0.999$} &
{\bfseries\boldmath $0.999$} & {\bfseries\boldmath $0.999$} & {\bfseries\boldmath $0.999$} &
{\bfseries\boldmath $0.999$} & {\bfseries\boldmath $0.99$} & {\bfseries\boldmath $0.999$} \\
\greyrule
{\bfseries\boldmath gradient clip} & \textbf{joint} & \textbf{joint} & \textbf{joint} &
\textbf{joint} & \textbf{joint} & \textbf{joint} & {\bfseries\boldmath per network} \\
\blackrule
configurations & $4$ & $4$ & $4$ & $4$ & $4$ & $4$ & $4$ \\
\greyrule
total, batch 1 & \multicolumn{7}{c}{$11$ units $\times$ $4$ weights $=44$ configurations,
$5{,}632$ copies} \\
\blackrule
\end{tabular}%
}
\caption{Batch 1, the seven algorithm ideas. Bold rows are swept. Under ``task reward used'',
$\sigma$ is the running standard deviation of the discounted sum of task rewards since the last
episode end. ``Joint'' clips the actor's, the critic's and the predictor's gradient together to
norm $0.5$; ``per network'' clips each to $0.5$ separately.}
\label{tab:pm-jax-batch1-alg-ideas}
\end{table}

\paragraph{Metrics and ``fixed''.} The metrics are those of
\Cref{sec:pm-jax-screening}, on each idea's own reward. An idea fixes the fall when its final-10\%
return is at least $90\%$ of its own peak and it still learns: an algorithm idea's final return
must exceed the control's peak minus one standard error; an environment idea, whose reward has a
different scale, must end with a share of episodes reaching the goal at least the control's share
at its peak.

\paragraph{Run status.} Submitted 2026-10-19 16:48 PT: eight units on RTX A4500 cards and three on
RTX 2080 Ti cards.
\studyunreadend % REPORT-STATUS

\end{document}
